% !TEX program = xelatex
\documentclass{英语试卷}
\begin{document}
\examtitle{宁波市 2024 学年第一学期高考模拟考试}{英语}
\par 宁波市 2024 学年第一学期高考模拟考试\par
\par 英语试卷\par
\begin{notice}
\par 1. 答卷前，考生务必将自己的姓名、准考证号填写在答题卡上。\par
\par 2. 回答选择题时，选出每小题答案后，用铅笔把答题卡上对应题目的答案标号涂黑。如需改动，用橡皮擦干净后，再选涂其他答案标号。回答非选择题时，将答案写在答题卡上，写在本试卷上无效。\par
\par 3. 考试结束后，将本试卷和答题卡一并交回。\par
\end{notice}
\parttitle{第一部分 听力（共两节，满分 30 分）}
\par 做题时，先将答案标在试卷上，录音内容结束后，你将有两分钟的时间将试卷上的答案转涂到答题卡上。\par
\sectitle{第一节（共5小题：每小题1.5分，满分7.5分）}
\par 听下面 5 段对话。每段对话后有一个小题，从题中所给的 A、B、C 三个选项中选出最佳选项。听完每段对话后，你都有 10 秒钟的时间来回答有关小题和阅读下一小题。每段对话仅读一遍。\par
\question{1. What did the woman ride?}
\begin{choices}(3)
  \item A horse.
  \item A cow.
  \item A camel.
\end{choices}
\question{2. What season is it now?}
\begin{choices}(3)
  \item Spring.
  \item Autumn.
  \item Winter.
\end{choices}
\question{3. What disturbed the woman during the night?}
\begin{choices}(3)
  \item A car starting.
  \item A dog barking.
  \item People arguing.
\end{choices}
\question{4. Why does the man want to leave the restaurant?}
\begin{choices}(3)
  \item He doesn't like the food.
  \item He can't stand the smell.
  \item He has been there too many times.
\end{choices}
\question{5. Where does the conversation probably take place?}
\begin{choices}(3)
  \item On a road.
  \item In a shop.
  \item At a concert.
\end{choices}
\sectitle{第二节（共15小题：每小题1.5分，满分22.5分）}
\par 听下面 5 段对话或独白。每段对话或独白后有几个小题，从题中所给的 A、B、C 三个选项中选出最佳选项。听每段对话或独白前，你将有时间阅读各个小题，每小题 5 秒钟；听完后，各小题将给出 5 秒钟的作答时间。每段对话或独白读两遍。 听下面一段对话，回答第 6、7 题。\par
\question{6. Why did the speakers go to France last year?}
\begin{choices}(3)
  \item They went on a study trip.
  \item They spent a vacation there.
  \item They visited some famous people.
\end{choices}
\question{7. How did the man find Picasso's paintings?}
\begin{choices}(3)
  \item Ugly.
  \item Realistic.
  \item Beautiful.
\end{choices}
\par 听下面一段对话，回答第 8 至 10 题。\par
\question{8. What is the relationship between the speakers?}
\begin{choices}(3)
  \item Father and daughter.
  \item Teacher and student.
  \item Boss and employee.
\end{choices}
\question{9. What is the girl doing with her phone?}
\begin{choices}(3)
  \item Studying.
  \item Playing games.
  \item Sending short messages.
\end{choices}
\question{10. What does the man concern most?}
\begin{choices}(3)
  \item The school work.
  \item The little pictures.
  \item The way of communication.
\end{choices}
\par 听下面一段对话，回答第11至13题。\par
\question{11. Where are the speakers?}
\begin{choices}(3)
  \item In Brazil.
  \item In China.
  \item In Argentina.
\end{choices}
\question{12. What is a characteristic of maté?}
\begin{choices}(3)
  \item Comforting.
  \item Refreshing.
  \item Calming.
\end{choices}
\question{13. What will the man do next?}
\begin{choices}(3)
  \item Sell the woman some tea.
  \item Drink tea with the woman.
  \item Show how the tea is produced.
\end{choices}
\par 听下面一段对话，回答第14至17题。\par
\question{14. What has been changed previously?}
\begin{choices}(3)
  \item More parking spaces.
  \item All new computers.
  \item More sports teams.
\end{choices}
\question{15. What does the woman intend to do next?}
\begin{choices}(3)
  \item Improve the food in the canteen.
  \item Better housing conditions on campus.
  \item Get more job opportunities for students.
\end{choices}
\question{16. What is the most important according to the woman?}
\begin{choices}(3)
  \item Enhancing the school's reputation.
  \item Increasing students' access to resources.
  \item Making life easier for parents and teachers.
\end{choices}
\question{17. What does the man long for?}
\begin{choices}(3)
  \item The school dance.
  \item The science fair.
  \item The spring concert.
\end{choices}
\par 听下面一段独白，回答第18至20题。\par
\question{18. When will the production of petrol cars stop?}
\begin{choices}(3)
  \item By 2050.
  \item By 2030.
  \item By 2025.
\end{choices}
\question{19. What does the speaker suggest as an alternative to current plastic bottles?}
\begin{choices}(3)
  \item Glass bottles.
  \item Metal bottles.
  \item Wood bottles.
\end{choices}
\question{20. What is the main idea of the talk?}
\begin{choices}(3)
  \item Dangers of global warming.
  \item New rules about electric cars.
  \item Ways to protect the environment.
\end{choices}
\parttitle{第二部分 阅读理解（共两节，满分50分）}
\sectitle{第一节（共15小题：每小题2.5分，满分37.5分）}
\par 阅读下列短文，从每题所给的 A、B、C 和 D 四个选项中，选出最佳选项。\par
\passagelabel{A}
\par Are you looking for your next great read? Whether it's for personal growth, inspiration, or an escape into an imaginative world, these newly released books have something for everyone.\par
\par The Power of Habit by Charles Duhigg\par
\par Charles Duhigg explores how habits shape our lives. Using real-life examples, he presents how understanding the "habit loop" can help improve health, productivity, and relationships. This book offers practical strategies for making lasting changes that lead to remarkable results.\par
\par 503 pages, 4.47 out of 5 stars\par
\par The Book of Joy by Emma Green and John Reynold\par
\par In this delightful book, Emma Green, a respected life coach, and John Reynold, a renowned psychologist, share insights into finding lasting happiness. Their engaging conversations, filled with wisdom and humor, offer a road map to living a joyful life despite challenges.\par
\par 291 pages, 4.11 out of 5 stars\par
\par Educated by Tara Westover\par
\par Educated is a powerful memoir that tells the story of Tara Westover, who was raised in a poor rural area and had no formal education until the age of 17. Her determination to educate herself and her journey to earning a PhD from Cambridge University are both inspiring and thought-provoking.\par
\par 464 pages, 4.19 out of 5 stars\par
\par The Midnight Library by Matt Haig\par
\par Nora Seed finds herself in a magical library where she has the chance to explore different lives she could have lived through each book, learning valuable lessons. This imaginative novel is perfect for readers who enjoy stories about second chances and self-discovery. 368 pages, 4.15 out of 5 stars\par
\question{21. Which of the following receives the highest rating?}
\begin{choices}(2)
  \item Educated.
  \item The Book of Joy.
  \item The Midnight Library.
  \item The Power of Habit.
\end{choices}
\question{22. Whose book appeals to readers in favor of lively discussions?}
\begin{choices}(2)
  \item Matt Haig's.
  \item Tara Westover's.
  \item Emma Green and John Reynold's.
  \item Charles Duhigg's.
\end{choices}
\question{23. What is the writing purpose of this text?}
\begin{choices}(2)
  \item To advertise some classic books.
  \item To express the author's preference.
  \item To recommend some quality books.
  \item To introduce different writing styles
\end{choices}
\passagelabel{B}
\par Crowded by defenders, Charlie Villanueva leaps toward the basket. One hand lifts the ball high, and for a moment he seems to hang in the air. Then he hammers the ball through the rim. Slam dunk! His career may look like a slam dunk, being a national champion at the University of Connecticut, and setting beginner records with the Toronto Raptors.\par
\par However, life wasn't always easy for him. His success is due to more than his ability to play basketball and he owes much of who he is today to a skin disease called alopecia areata, which is not infectious but is caused by an error in the immune system that triggers germ-fighting cells to attack the body's own hair-growing cells. Charlie was affected at age 10 and lost all his hair by age 12. This made him a target for bullying, leading him to withdraw and hide under caps and hoods\par
\par With the support of his mother, Charlie found refuge in basketball. He channeled his frustrations and stress into motivation, improving his skills to the point where his talent overshadowed his condition. The bullying stopped, and his basketball career took off.\par
\par All along the way, Charlie has consistently dedicated himself to helping others. Ever since his first year as a new player, he has served as a spokesperson for the National Alopecia Areata Foundation, and went on to establish the Charlie Villanueva Foundation, an organization aimed at supporting families with children suffering from alopecia and stopping all forms of bullying. As the son of immigrants from the Dominican Republic, he has even worked towards ending poverty in that country.\par
\par Charlie inspires young people to embrace their individuality and have confidence in themselves. "We're all unique in our own ways," he says, encouraging them to remember that with dedication and self-belief, any dream is within reach.\par
\question{24. Why is Charlie's slam dunk mentioned in Paragraph 1?}
\begin{choices}(2)
  \item To compare him with other champions.
  \item To highlight his skills and achievements.
  \item To explain how a slam dunk is performed.
  \item To demonstrate a specific moment in a game.
\end{choices}
\question{25. What made Charlie avoid his peers?}
\begin{choices}(2)
  \item An infectious skin disease.
  \item His being made fun of.
  \item An error in the immune system.
  \item His average ability to play basketball.
\end{choices}
\question{26. Which of the following best describes Charlie?}
\begin{choices}(2)
  \item Humble and responsible.
  \item Strong-willed and helpful.
  \item Innovative and generous.
  \item Good-natured and tolerant.
\end{choices}
\question{27. What message does the text convey?}
\begin{choices}(2)
  \item Fortune favors the prepared.
  \item Actions speak louder than words.
  \item What doesn't kill you makes you stronger.
  \item A journey of a thousand miles begins with a single step.
\end{choices}
\passagelabel{C}
\par Among bearing a reputation for various characteristics ranging from the good to the bad, millennials (千禧一代) have notably been labeled the "job-hopping" generation. A Deloitte survey revealed that 49\% of millennials might quit their jobs within the next two years. While some from older generations may interpret this as a sign of laziness, a deeper look suggests more complex motivations.\par
\par At first glance, there are a variety of reasons and circumstances that may lead a person to job hop occasionally. However, beyond the conventional reasons like obtaining better salaries or greater job stability, many millennials have been increasingly motivated by value-driven goals. Surveys show that some go to the extent of suffering pay cuts to pursue positions that provide better learning opportunities. Additionally, some even quit their positions for companies which are aligned with their value system better. In response to such a trend, employers are taking steps to improve workplace culture and to pay more attention to the professional growth and development of individuals.\par
\par Visibly, Millennials' boldness in pursuing new job opportunities is also shaped by emerging economic trends and novel business models, such as the gig economy, where self-employed and temporary work caters to those who value flexibility and work-life balance. Additionally, social media has created new career paths, from blogging to content creation, leading more millennials to leave traditional desk jobs for unconventional careers. Hence, one can't say for certain that job-hopping indicates laziness. Instead, it shows millennials' wandering off the beaten track to explore new opportunities and development.\par
\par However, job-hopping remains controversial. Some employers appreciate the diversity in experience and skills job-hopping brings about, while others view it negatively, seeing it as a lack of commitment. There is concern that millennials, raised in relatively prosperous times, may possess less patience and resilience. Yet, amidst a materialistic society, millennial generation's pursuit of meaning in their careers offers inspiration. They should never be discouraged to define new horizons of success when their career decisions are driven by healthy ambition and an appetite for upskilling.\par
\question{28. What phenomenon does the first paragraph describe?}
\begin{choices}(2)
  \item Millennials switch jobs frequently.
  \item Millennials are reluctant to seek jobs.
  \item Millennials have fewer job opportunities.
  \item Millennials are particular about their jobs.
\end{choices}
\question{29. The underlined phrase "are aligned with" in Paragraph 2 can be replaced by \blanknum{}.}
\begin{choices}(2)
  \item go beyond
  \item correspond to
  \item conflict with
  \item compete against
\end{choices}
\question{30. The emerging gig economy attracts Millennials in that it allows for \blanknum{}.}
\begin{choices}(2)
  \item decent income
  \item high social status
  \item definite freedom
  \item self-paced schedule
\end{choices}
\question{31. What is the author's view on the job-hopping generation?}
\begin{choices}(2)
  \item They are characterized by laziness.
  \item They lack diverse work experience.
  \item They prioritize individual fulfillment.
  \item They should bear more responsibility.
\end{choices}
\passagelabel{D}
\par For thousands of years, the Taklamakan Desert preserved more than the bodies of those buried in China's Tarim Basin. Archaeologists found pebble-sized pieces of a yellowish substance: the world's oldest cheese, scattering around the necks of some bodies in Xiaohe, a Bronze Age cemetery. A recent study, published in Cell, analyzes the DNA of this ancient dairy, shedding light on how it was made and how its production spread.\par
\par In a previous study, researchers took tiny samples of the 3,500-year-old cheese attached to the mummies' necks. An analysis of proteins in those bits unveiled the presence of Lactobacillus kefiranofaciens, a microbe(发酵菌) used to produce a type of fermented cheese called kefir.\par
\par "That was really amazing", says Qiaomei Fu, a scientist at Chinese Academy of Sciences in Beijing. Having worked with ancient human DNA, she wanted to study the microbe’s DNA. However, extracting complete DNA from such ancient samples proved challenging, as it broke down over time and mixed with genetic information from the environment. In 2014, Fu began designing special molecules (分子) that could efficiently capture the microbe's ancient DNA from the cheese. After years of work, Fu and her team managed to reconstruct 92 percent of the ancient microbe’s DNA.\par
\par From this microbe's DNA, the team could explore the history of the cheese's production. Kefir is made by mixing milk with an existing microbe, and as it spread, so did the microbes used to make it. By comparing the DNA from the Xiaohe cheese with that of others, two branches emerged: one from Europe and the Pacific Islands, and the other from Tibet and East Asia, which suggests a possible route for kefir's spread from Xinjiang to East Asia, possibly through migration.\par
\par The DNA analysis also pointed to the interaction of ancient peoples. It uncovered that samples from cows and goats were found. Interestingly, the goat DNA resembled that of other ancient Central Asian goats, suggesting the mixing of population in Xinjiang with other Bronze Age populations. Researchers anticipate that further analysis of microbes in the cheese, along with other biological samples, could provide deeper insights into interactions between ancient populations.\par
\question{32. What's the significance of the DNA analysis of the ancient cheese?}
\begin{choices}(2)
  \item It reflects the civilization of the Bronze Age
  \item It proves China to be the first to produce dairy products.
  \item It uncovers the mysterious buried bodies in Xiaohe cemetery.
  \item It reveals production methods and transmission routes of the cheese.
\end{choices}
\question{33. What obstacle did Fu's team encounter in their study?}
\begin{choices}(2)
  \item The complete set of DNA was hard to obtain.
  \item The team lacked working experience with DNA.
  \item The previous study proved the absence of a microbe.
  \item The cheese samples were resistant to DNA extraction methods.
\end{choices}
\question{34. What can be inferred from the research findings?}
\begin{choices}(2)
  \item Ancient people migrated typically with cows and goats.
  \item Cheese played an essential role in ancient people's diets.
  \item Kefir originated in Europe and spread to Tibet and East Asia.
  \item Kefir was possibly exchanged between peoples across different regions.
\end{choices}
\question{35. What can be a suitable title for the text?}
\begin{choices}(2)
  \item What's in the Oldest Cheese?
  \item Can DNA Tell Secrets of Ancient Diets?
  \item Did Ancient Civilizations Perfect Cheese-Making?
  \item How Did the Oldest Cheese Connect with Cultural Exchange?
\end{choices}
\sectitle{第二节（共5小题；每小题2.5分，满分12.5分）}
\par 阅读下面短文，从短文后的选项中选出能填入空白处的最佳选项。选项中有两项为多余选项，\par
\par Virtual internships (实习) are now a common occurrence as we adjust to the "new normal" when it comes to working. \blanknum{36} Still, there are many positives that can come from virtual internships that you should consider. Use it for resume building\par
\par \blanknum{37} It showcases your ability to manage your time effectively and handle tasks on your own. Besides, securing glowing recommendations from a remote employer highlights skills you gained through remote work.\par
\par Having a job can cost quite a bit when you take the commute (通勤) into consideration. With remote work, the costs are significantly lower, if they exist at all. \blanknum{38} While being in person has a certain feel difficult to let go of, working from home is arguably one of the best ways to store up some money\par
\par Make the most of your time savings\par
\par Additionally, without that commute, you gain some extra time to get things done that you otherwise wouldn't have been able to. This includes many things, such as watching your favorite shows and cooking homemade meals on a nightly basis. \blanknum{39}\par
\par Treat it as experience\par
\par Lastly, keep in mind that a virtual internship is still an internship. Experience is experience, no matter how you obtain it. \blanknum{40} For that reason, understand that. a virtual internship is still experience, and that it's still a feasible option to include on your resume.\par
\begin{choices}(1)
  \item That being said, it can be easily viewed as a negative.
  \item You get to work from comfort home, no commute required.
  \item It is the way of the future, at least, for the foreseeable future.
  \item Gas, parking fees, and even grabbing lunch can quickly add up.
  \item It doesn't seem like much, but the extra time can be a real game-changer.
  \item While it's easy to dismiss a virtual internship, you'd be missing an opportunity
  \item As for building a professional profile, remote work can really sell your experience.
\end{choices}
\parttitle{第三部分 语言运用（共两节，满分 30 分）}
\sectitle{第一节（共15小题；每小题1分，满分15分）}
\par 阅读下面短文，从短文后各题所给的 A、B、C 和 D 四个选项中，选出可以填入空白处的最佳选项，\par
\par The first day of school our professor introduced someone we didn't know. I looked around to find a (an) \blanknum{41} little lady beaming up at me with a smile that lit up the room.\par
\par Her name was Rose and she was eighty-seven years old. She \blanknum{42} asked me for a giant hug, and \blanknum{43} her seeking adventure and love in college. Actually, she had a \blanknum{44} purpose— fulfilling her dream\par
\par of a college education.\par
\par At the end of the semester we invited her to speak at our football banquet. She \blanknum{45} her throat and began, "There are some \blanknum{46} if you want to stay young, be happy and achieve success. You have to laugh and find humor each day. You've got to have dreams and no \blanknum{47}. The only people who \blanknum{48} death are those with pity and disappointment."\par
\par Eventually, she \blanknum{49} her speech by courageously singing "The Rose". She challenged each of us to study the \blanknum{50} and live them out in our daily lives.\par
\par At year's end, Rose finished the college \blanknum{51} she had begun all those years ago. One week after graduation Rose passed away \blanknum{52} in her sleep. Over two thousand college students attended her funeral to \blanknum{53} the wonderful woman who taught \blanknum{54} that it's never too late to be all you can \blanknum{55} be.\par
\cz{41}{exhausted}{wrinkled}{shaking}{demanding}
\cz{42}{enthusiastically}{sympathetically}{cautiously}{determinedly}
\cz{43}{complained about}{boasted about}{joked about}{dreamed about}
\cz{44}{newly-found}{widely-accepted}{long-held}{well-designed}
\cz{45}{examined}{cleared}{tightened}{scratched}
\cz{46}{messages}{lines}{reasons}{secrets}
\cz{47}{failures}{faults}{troubles}{regrets}
\cz{48}{fear}{defeat}{escape}{face}
\cz{49}{delivered}{concluded}{prepared}{repeated}
\cz{50}{lyrics}{tunes}{rhythms}{melodies}
\cz{51}{struggle}{application}{degree}{system}
\cz{52}{dramatically}{tragically}{peacefully}{courageously}
\cz{53}{miss}{mourn}{remember}{praise}
\cz{54}{on purpose}{by chance}{by example}{on schedule}
\cz{55}{possibly}{occasionally}{noticeably}{naturally}
\sectitle{第二节（共10小题；每小题1.5分，满分15分）}
\par 阅读下面短文，在空白处填入 1 个适当的单词或括号内单词的正确形式。\par
\par The "Italian Contemporary Art and Artists on the Silk Road" exhibition, \blanknum{56} (hold) at the Grand Tang Mall in Xi'an, drew over \blanknum{60},000 visitors during its 37-day run. Described by art \blanknum{57} (enthusiastic) Zhang Yin, this unique event displayed the integration of Chinese tradition with Western modernity, offering an \blanknum{58} (appeal) journey transcending time and space.\par
\par The exhibition's route began in Italy and spanned countries along the Silk Road, ending in Xi'an, \blanknum{59} route's origin. Pasquale Terracciano, from Italy's Ministry of Foreign Affairs, viewed the exhibition \blanknum{60} a symbol of renewal and dialogue, highlighting the role of art in bridging cultural divides.\par
\par \blanknum{61} (historic), the Silk Road promoted cultural exchange between China and Italy, which\blanknum{63} (evidence) by Marco Polo's travels and the spread of Chinese culture to Italy. The Xi' an Art Museum, together with Italian peers, preserves this legacy, having previously hosted the "Beyond the Wall: Xi' an Contemporary Art Exhibition" in Florence.\par
\par \blanknum{63} (feature) 31 pieces from 33 Italian artists, the exhibition presented a diverse range of creative expressions, reflecting Italy's artistic evolution from the 1960s to today. Notable was "What holds it all together," a chinoiserie-inspired piece made of porcelain and organic glass, \blanknum{64} exemplified the enduring cultural bond between the two nations.\par
\par Though the exhibition closed its doors in Xi'an, the echo of its profound impact on the cultural landscape \blanknum{65} (continue) to resonate, fostering a new era of creative harmony between China and Italy.\par
\parttitle{第四部分 写作（共两节，满分 40 分）}
\sectitle{第一节（满分15分）}
\par 学校英语俱乐部将组织关于“生活中具有重大影响的科技产品”的讨论，请你作为代表发言，内容包括：\par
\par （1）介绍一项产品及其影响；\par
\par （2）对该产品的展望。注意：\par
\par (1) 写作词数应为 80 左右;\par
\par (2) 请按如下格式在答题纸的相应位置作答。\par
\par Good morning, everyone!\par
\sectitle{第二节（满分25分）}
\par 阅读下面材料，根据其内容和所给段落开头语续写两段，使之构成一篇完整的短文。\par
\par The sun was setting over Springdale, casting shadows across the busy community center. A group of teenagers, the members of the Springdale Youth Committee, gathered at a table. They looked determined and worried. Among them was Emily, a smart but often unnoticed girl who loves helping others.\par
\par For weeks, they had been organizing a fundraising event for the local homeless shelter, as the shelter struggled to assist the increasing number of individuals seeking help. The big day was one week away, and everything was set — until they discovered another group had a similar charity event on the same day. This could divide the people who might come and the donations they hoped to get.\par
\par "We can't change the date now," said Tom, a senior member, his voice filled with frustration. "We've already sent out invitations, and the place is booked." Emily listened quietly to the discussion. She knew changing the date wasn't possible, but she also realized competing with another worthy cause could hurt both sides.\par
\par After the meeting ended without a solution, Emily decided to take action. She spent the next day looking into ways to work with the other group, hoping to find a way that would help both the homeless shelter and the other cause.\par
\par The next day, she came to the committee with a suggestion. "What if we join our events?" she asked. "We could share resources, work together on promoting, and most importantly, make sure both causes are supported by the community."\par
\par The members looked at each other, some doubtful, others interested. "That's a great idea, Emily, but how will we get them to agree?" Laura, another member, asked. Emily had already thought this through. "I'll talk to their organizer and show them the benefits of a joint event." She knew the path ahead would be challenging, but she was determined to make it happen for Springdale and its people. 注意：\par
\par (1) 续写词数应为 150 左右;\par
\par (2) 请按如下格式在答题卡的相应位置作答。\par
\writingpara{Paragraph 1: The committee agreed to let Emily lead this effort.}
\writinglines{8}
\writingpara{Paragraph 2: After days of negotiation, Emily finally got a positive reply.}
\writinglines{8}
\end{document}
